\section{Introduction}

A pipeline that produces more numbers than its authors can trace by hand has to
check itself, and one standard way to do it is redundancy: compute every
quantity that will be published by two routes built on different technology,
compare them, and refuse to exit successfully unless they agree. The check is
cheap, mechanical, and needs no ground truth, which is why it appeals wherever
findings are generated faster than they can be audited, including in systems
that propose and test hypotheses automatically \citep{lu2024aiscientist}. The
same intuition, that agreement between routes stands in for a correctness
signal, underwrites self-consistency decoding \citep{wang2023selfconsistency}
and the use of agreement between generated code and generated tests to rank
programs \citep{chen2023codet}.

This paper is a case study of one such gate failing, in a pipeline we built. It
is a single incident, and we do not claim that it tells us how often redundancy
gates fail. What it can do is show one failure completely: the gate, the wrong
and corrected numbers, the definition that caused the error, and what did and
did not catch it.

The pipeline compares two independently maintained public catalogues of
artificial Earth-orbiting objects and emits an RDF graph of typed defect
assertions. A dual-computation gate counts each defect class twice, once
set-based in Python from the source files and once by SPARQL over the emitted
graph, and exits non-zero on any disagreement. On the run whose numbers were
published it printed \texttt{ALL CROSS-CHECKS AGREE}. Three of the seven counts
it compared were wrong, and the largest was overstated more than fourfold. Both
paths imported the same classification constants, and those constants encoded a
misreading of the source's vocabulary of flight-phase status codes. Two correct
implementations of one wrong reading agree perfectly, so the error was
common-mode and the gate was silent by construction. The multiversion
programming literature established long ago that separately developed versions
of a program do not fail independently
\citep{knight1986experimental, eckhardt1985theoretical, littlewood1989conceptual};
our case is a data-pipeline instance of that result in which the shared artifact
is not a specification but an executable set of constants.

The contributions are four, each bounded to this pipeline.
First, the mechanism, with an object-level ledger that reconciles every figure
across versions of the count (Sections~\ref{sec:failure}
and~\ref{sec:mechanism}, Appendix~\ref{app:ledger}).
Second, three checks that go back to the source's own documentation, each
implemented and run on the defective code and on its correction
(Sections~\ref{sec:caught} and~\ref{sec:gates}).
Third, a second-order failure. Checking our correction against every object's
phase history showed that it misread the same vocabulary in a different way, and that
none of our three checks had noticed (Section~\ref{sec:second}).
Fourth, an experiment on whether a code-generating model asked for an
independent verification path produces one (Section~\ref{sec:automation}).

Redundancy remains worth having; it verifies implementation. The errors that
reached publication here were errors of meaning, which two paths built on shared
definitions cannot test.
